\documentclass[a4paper]{article}

% Español (México) con XeLaTeX: fontspec para fuentes Unicode, polyglossia
% para la división silábica y los nombres del documento en la variante mexicana.
\usepackage{fontspec}
\usepackage{polyglossia}
\setdefaultlanguage[variant=mexican]{spanish}
% Los nombres chinos van como «pinyin (original)»: xeCJK da a los caracteres
% Han su propia fuente; sin ella XeLaTeX los omite con un aviso.
% FandolSong no cubre algunos caracteres de nombres (U+73FA); WenQuanYi Zen Hei sí.
\usepackage[AutoFallBack=true]{xeCJK}
\setCJKmainfont{FandolSong-Regular.otf}
\setCJKfallbackfamilyfont{\CJKrmdefault}{WenQuanYi Zen Hei}
% polyglossia no traduce los nombres de listings.
\AtBeginDocument{\providecommand{\lstlistingname}{}\renewcommand{\lstlistingname}{Listado}%
  \providecommand{\lstlistlistingname}{}\renewcommand{\lstlistlistingname}{Índice de listados}}
\usepackage{amsmath,amssymb}
\usepackage{graphicx}
\usepackage[margin=2.5cm]{geometry}
\usepackage[most]{tcolorbox}
\usepackage{etoolbox}
\usepackage{subcaption}
\usepackage{listings}
\usepackage{hyperref}
\usepackage{booktabs}
\usepackage{float}

\newtcolorbox{knowledgebox}[1]{
    enhanced, colback=blue!5!white, colframe=blue!75!black, colbacktitle=blue!75!black,
    coltitle=white, fonttitle=\bfseries, title={#1}, boxrule=1pt, sharp corners
}

\newtcolorbox{importantbox}[1]{
    enhanced, colback=yellow!10!white, colframe=yellow!80!black, colbacktitle=yellow!80!black,
    coltitle=black, fonttitle=\bfseries, title={#1}, sharp corners
}

\newtcolorbox{warningbox}[1]{
    enhanced, colback=red!5!white, colframe=red!75!black, colbacktitle=red!75!black,
    coltitle=white, fonttitle=\bfseries, title={#1}, sharp corners
}

\lstset{
    basicstyle=\ttfamily\small,
    keywordstyle=\color{blue},
    stringstyle=\color{red!60!black},
    commentstyle=\color{green!60!black},
    breaklines=true,
    frame=single,
    numbers=left,
    numberstyle=\tiny\color{gray},
    captionpos=b,
    extendedchars=true
}

\newcommand{\notetitle}{CS224R Lecture 17: Avance de la inteligencia robótica mediante aprendizaje por refuerzo}
\newcommand{\noteauthors}{Elaboradas a partir de materiales públicos del curso}
\newcommand{\notedate}{Primavera de 2025}
\newcommand{\videochannel}{Stanford Online}
\newcommand{\videopublishdate}{2025}
\newcommand{\videoduration}{Aproximadamente 60 minutos}
\newcommand{\videourl}{https://cs224r.stanford.edu/}
\newcommand{\videocoverpath}{cover.jpg}
\newcommand{\repourl}{https://github.com/hqhq1025/ai-course-notes}


% Footer with repo info
\usepackage{fancyhdr}
\pagestyle{fancy}
\fancyhf{}
\fancyfoot[C]{\thepage}
\fancyfoot[R]{\footnotesize\color{gray}\href{\repourl}{AI Course Notes}}
\renewcommand{\headrulewidth}{0pt}
\renewcommand{\footrulewidth}{0pt}

\begin{document}
\begin{titlepage}
\centering
{\Large Notas del curso\par}
\vspace{1.2cm}
{\huge\bfseries \notetitle\par}
\vspace{0.8cm}
{\large \noteauthors\par}
\vspace{0.3cm}
{\large \notedate\par}
\vspace{1.1cm}

\ifdefempty{\videocoverpath}{
{\small No se proporcionó la imagen de portada.\par}
}{
\includegraphics[width=0.82\textwidth,height=0.45\textheight,keepaspectratio]{\videocoverpath}\par
}

\vfill
\begin{tcolorbox}[width=0.9\textwidth, colback=black!2!white, colframe=black!60, sharp corners]
\textbf{Autor o canal del video}: \videochannel\par
\textbf{Fecha de publicación}: \videopublishdate\par
\textbf{Duración del video}: \videoduration\par
\textbf{Ponente}: Ashish Kumar (Guest Lecture)\par
\textbf{Página del curso}: \href{https://cs224r.stanford.edu/}{cs224r.stanford.edu}\par
\textbf{Página del proyecto}: \href{\repourl}{\nolinkurl{\repourl}}\par
\end{tcolorbox}
\end{titlepage}

\tableofcontents
\newpage

%% --- Inicio del contenido --- %%

\section{Aprendizaje por imitación vs. aprendizaje por refuerzo: comparación práctica}

Esta clase la imparte como invitado Ashish Kumar y se centra en las aplicaciones prácticas del RL en robótica, en particular en locomoción y manipulación. El ponente comienza comparando el aprendizaje por imitación y el aprendizaje por refuerzo «que funcionan en la práctica».

\begin{importantbox}{Rasgos prácticos de IL vs. RL}
\begin{itemize}
    \item \textbf{Aprendizaje por imitación}:
    \begin{itemize}
        \item Es eficiente en datos y conviene en escenarios donde se pueden obtener demostraciones de alta calidad
        \item Su rendimiento máximo está limitado por la capacidad del demostrador
        \item Difícilmente supera el nivel humano
    \end{itemize}
    \item \textbf{Aprendizaje por refuerzo}:
    \begin{itemize}
        \item Puede descubrir políticas que superan a las humanas
        \item Suele requerir un entorno de simulación (sim-to-real transfer)
        \item El diseño de la recompensa es el principal desafío
        \item El entrenamiento es largo, pero la inferencia es eficiente
    \end{itemize}
\end{itemize}
\end{importantbox}

\subsection{Cuándo elegir RL}

\begin{knowledgebox}{Condiciones en que conviene el RL}
El RL tiene ventaja sobre el IL en los siguientes escenarios:
\begin{itemize}
    \item Se dispone de un \textbf{simulador} de alta calidad
    \item El objetivo de la tarea puede cuantificarse con una función de recompensa
    \item Se necesita \textbf{superar} el nivel de los demostradores humanos
    \item Los datos de demostración son difíciles de obtener (por ejemplo, en operaciones peligrosas)
    \item Se requiere \textbf{robustez} ante cambios en el entorno
\end{itemize}
\end{knowledgebox}

\subsection{Resumen de la sección}

IL y RL tienen cada uno sus escenarios de aplicación, y en los sistemas reales suelen combinarse (por ejemplo, IL warm-start + RL fine-tuning).

\section{Sim-to-Real Transfer}

\subsection{Por qué se necesita la simulación}

Hacer RL en robots reales plantea grandes desafíos:
\begin{itemize}
    \item La recolección de datos es lenta (se ejecuta en tiempo real)
    \item Riesgos de seguridad (el robot puede dañarse a sí mismo o dañar el entorno)
    \item Es difícil paralelizar a gran escala
\end{itemize}

Los entornos de simulación resuelven estos problemas, pero introducen la \textbf{sim-to-real gap}: la diferencia entre la simulación y el mundo real.

\begin{importantbox}{Orígenes de la Sim-to-Real Gap}
\begin{itemize}
    \item \textbf{Diferencias de dinámica}: el simulador no puede modelar a la perfección la física real (fricción, contacto, deformación)
    \item \textbf{Diferencias de percepción}: la visión y el tacto simulados difieren de los sensores reales
    \item \textbf{Diferencias de actuadores}: los motores reales tienen retardo, ruido y características no lineales
    \item \textbf{Diferencias de entorno}: el entorno real tiene objetos no modelados, cambios de iluminación, etc.
\end{itemize}
\end{importantbox}

\begin{figure}[H]
\centering
\includegraphics[width=0.9\textwidth]{slides-images/slide-12.jpg}
\caption{El ponente explica con un marco de entrenamiento con codificador de parámetros extrínsecos que aleatorizar en simulación factores como la masa, la fricción y el terreno es un paso de preparación esencial antes de transferir la política a un robot real.}
\end{figure}

\subsection{Domain Randomization}

\begin{importantbox}{Aleatorización de dominio}
Domain Randomization es uno de los métodos principales para cerrar la sim-to-real gap:

En la simulación se aleatorizan diversos parámetros del entorno (coeficientes de fricción, tamaño de los objetos, masa, apariencia visual, ruido de los sensores, etc.), de modo que la política, ya en la fase de entrenamiento, «haya visto» múltiples variantes posibles del entorno.

Supuesto central: si la política se desempeña bien en entornos simulados suficientemente diversos, el mundo real puede considerarse una de esas variantes.
\end{importantbox}

\begin{warningbox}{Limitaciones de Domain Randomization}
\begin{itemize}
    \item Una aleatorización excesiva puede volver la política demasiado conservadora
    \item Es necesario asegurar que los parámetros del mundo real estén dentro del rango de aleatorización
    \item Algunos fenómenos físicos (como el contacto con cuerpos blandos) son difíciles de cubrir incluso con aleatorización
    \item Elegir qué parámetros aleatorizar y en qué rango requiere conocimiento del dominio
\end{itemize}
\end{warningbox}

\subsection{Resumen de la sección}

La sim-to-real transfer es el puente fundamental para aplicar RL a robots reales, y la domain randomization es actualmente el método más utilizado.

\section{Locomoción robótica: locomoción con patas}

\subsection{Caminata de robots cuadrúpedos}

El ponente presentó casos en los que se usa RL para entrenar robots cuadrúpedos (como Unitree Go1/A1) para caminar sobre distintos terrenos.

\begin{importantbox}{Marco de RL para la locomoción con patas}
\begin{itemize}
    \item \textbf{Estado}: ángulos de las articulaciones, velocidades angulares, datos de la IMU (ángulo de inclinación, velocidad angular)
    \item \textbf{Acción}: ángulo objetivo de cada articulación (seguido por un controlador PD)
    \item \textbf{Recompensa}: velocidad de avance + penalización por energía + penalización por postura + recompensa por frecuencia de pasos
    \item \textbf{Entrenamiento}: se realiza en simuladores paralelos en GPU como Isaac Gym, que pueden ejecutar miles de entornos al mismo tiempo
\end{itemize}
\end{importantbox}

\subsection{Marco Teacher-Student}

\begin{knowledgebox}{Destilación de información privilegiada}
Un paradigma de entrenamiento sim-to-real muy eficaz:
\begin{enumerate}
    \item \textbf{Fase Teacher}: se entrena en simulación una política teacher a la que se le permite acceder a \textbf{información privilegiada} (como el mapa de alturas exacto del terreno, la masa de los objetos y el coeficiente de fricción)
    \item \textbf{Fase Student}: se entrena una política student que solo usa \textbf{sensores disponibles en el mundo real} (como la propiocepción y las cámaras) y que imita el comportamiento del teacher mediante destilación
\end{enumerate}
El teacher puede aprender mejor (porque dispone de más información), mientras que el student aprende a inferir información equivalente a partir de sensores limitados.
\end{knowledgebox}

\begin{figure}[H]
\centering
\includegraphics[width=0.9\textwidth]{slides-images/slide-16.jpg}
\caption{Lo esencial en la fase de despliegue no es «reutilizar directamente la política de simulación», sino estimar en línea los parámetros extrínsecos a partir del historial de observaciones y corregir el controlador de forma continua; esta es también la fuente principal de la robustez de los robots con patas modernos.}
\end{figure}

\subsection{Resumen de la sección}

RL + sim-to-real se ha convertido en el paradigma dominante para el control de robots con patas, y el marco teacher-student resuelve de forma eficaz el problema de la percepción limitada.
\section{Manipulación robótica}

\subsection{Los retos de la manipulación diestra}

La manipulación diestra (por ejemplo, sujetar y manipular objetos con una mano de varios dedos) es más difícil que la locomoción:
\begin{itemize}
    \item La dinámica de contacto es fuertemente no lineal
    \item El espacio de estados tiene más dimensiones
    \item Las tareas son más diversas
    \item La Sim-to-real gap es más grave en el modelado del contacto
\end{itemize}

\subsection{Aplicaciones del RL en la manipulación}

El ponente presentó varios casos de manipulación diestra con RL, entre ellos:
\begin{itemize}
    \item Girar un cubo de Rubik (OpenAI Rubik's Cube)
    \item Hacer girar un bolígrafo (pen spinning)
    \item Sujeción diestra (dexterous grasping)
\end{itemize}

Todas estas tareas siguen el mismo enfoque: simulación paralela a gran escala, domain randomization y teacher-student.

\subsection{Resumen de la sección}

El RL ya ha mostrado una capacidad superior a la humana en diversas tareas de manipulación robótica, pero la calidad del simulador sigue siendo un factor limitante decisivo.

\section{Robots humanoides}

\subsection{Marcha bípeda y control de cuerpo completo}

El ponente presentó brevemente los avances de RL aplicado a robots humanoides:
\begin{itemize}
    \item Marcha bípeda estable
    \item Coordinación de cuerpo completo (caminar y manipular al mismo tiempo)
    \item Robustez ante perturbaciones de fuerzas externas
\end{itemize}

\begin{knowledgebox}{Retos particulares de los robots humanoides}
En comparación con los robots cuadrúpedos, el entrenamiento con RL de robots humanoides enfrenta retos mayores:
\begin{itemize}
    \item Apoyo bípedo inestable (margen de estabilidad más estrecho)
    \item Más grados de libertad (más de 30 articulaciones en todo el cuerpo)
    \item Las extremidades superiores e inferiores deben coordinarse
    \item Las caídas tienen consecuencias más graves (el hardware es costoso)
\end{itemize}
\end{knowledgebox}

\subsection{Resumen de la sección}

Los robots humanoides representan el reto de frontera de RL for robotics; los avances actuales son rápidos, pero todavía falta para lograr un despliegue estable.
\section{Síntesis y ampliación}

\begin{enumerate}
    \item RL e IL tienen ventajas distintas; los sistemas reales suelen combinarlos
    \item La sim-to-real transfer es el flujo de trabajo central para aplicar RL a robots reales
    \item La domain randomization y el marco teacher-student son técnicas fundamentales para cerrar la sim-to-real gap
    \item RL ya logró avances decisivos en locomoción con patas y en manipulación diestra
    \item Los robots humanoides representan el reto de la siguiente generación
    \item La calidad de los simuladores y el diseño de recompensas siguen siendo el cuello de botella de las aplicaciones reales
\end{enumerate}

\begin{knowledgebox}{Lecturas adicionales}
\begin{itemize}
    \item Kumar et al., ``Rapid Motor Adaptation (RMA),'' RSS 2021
    \item Tobin et al., ``Domain Randomization for Transferring Deep Neural Networks from Simulation to the Real World,'' IROS 2017
    \item OpenAI, ``Solving Rubik's Cube with a Robot Hand,'' 2019
    \item Rudin et al., ``Learning to Walk in Minutes Using Massively Parallel Deep RL,'' CoRL 2022
    \item Radosavovic et al., ``Real-World Humanoid Locomotion with RL,'' Science Robotics 2024
\end{itemize}
\end{knowledgebox}

\end{document}
